\textbf{Phần 1. Trắc nghiệm lựa chọn}
\setcounterpageref{ex}{0}

\begin{ex} %[1D10V40]
Hình tam giác có bao nhiêu cạnh?
\choice
{$2$}
{$4$}
{$1$}
{\True $3$}
\loigiai{
Hình tam giác có đúng 3 cạnh.
\textbf{Đáp số:} $3$.\par
\textbf{Đáp án của câu hỏi là D}
}
\end{ex}

\begin{ex} %[1D10V40]
Hình chữ nhật có bao nhiêu góc vuông?
\choice
{$3$}
{\True $4$}
{$5$}
{$2$}
\loigiai{
Hình chữ nhật có đúng 4 góc vuông.
\textbf{Đáp số:} $4$.\par
\textbf{Đáp án của câu hỏi là B}
}
\end{ex}

\begin{ex} %[1D10V40]
Khối lập phương có bao nhiêu mặt?
\choice
{$5$}
{$7$}
{\True $6$}
{$4$}
\loigiai{
Khối lập phương có đúng 6 mặt phẳng đều là hình vuông.
\textbf{Đáp số:} $6$.\par
\textbf{Đáp án của câu hỏi là C}
}
\end{ex}

\begin{ex} %[1D10V40]
Khối hộp chữ nhật có bao nhiêu đỉnh?
\choice
{\True $8$}
{$7$}
{$9$}
{$6$}
\loigiai{
Khối hộp chữ nhật có đúng 8 đỉnh.
\textbf{Đáp số:} $8$.\par
\textbf{Đáp án của câu hỏi là A}
}
\end{ex}

\begin{ex} %[1D10V40]
Đồng hồ có kim ngắn chỉ số 8, kim dài chỉ số 12. Hỏi đồng hồ chỉ mấy giờ?
\choice
{$7\text{ giờ}$}
{$9\text{ giờ}$}
{$6\text{ giờ}$}
{\True $8\text{ giờ}$}
\loigiai{
Đồng hồ chỉ đúng 8 giờ.
\textbf{Đáp số:} $8$.\par
\textbf{Đáp án của câu hỏi là D}
}
\end{ex}

\begin{ex} %[1D10V40]
Một tuần lễ có bao nhiêu ngày?
\choice
{$6\text{ ngày}$}
{\True $7\text{ ngày}$}
{$8\text{ ngày}$}
{$5\text{ ngày}$}
\loigiai{
Một tuần lễ có đúng 7 ngày.
\textbf{Đáp số:} $7$.\par
\textbf{Đáp án của câu hỏi là B}
}
\end{ex}

\begin{ex} %[1D10V40]
Tính: $35\text{ cm} + 24\text{ cm}$. Kết quả là bao nhiêu xăng-ti-mét?
\choice
{$58\text{ cm}$}
{$60\text{ cm}$}
{\True $59\text{ cm}$}
{$57\text{ cm}$}
\loigiai{
Ta có: $35 + 24 = 59$ (cm).
\textbf{Đáp số:} $59$.\par
\textbf{Đáp án của câu hỏi là C}
}
\end{ex}

\begin{ex} %[1D10V40]
Tính: $78\text{ cm} - 35\text{ cm}$. Kết quả là bao nhiêu xăng-ti-mét?
\choice
{\True $43\text{ cm}$}
{$42\text{ cm}$}
{$44\text{ cm}$}
{$41\text{ cm}$}
\loigiai{
Ta có: $78 - 35 = 43$ (cm).
\textbf{Đáp số:} $43$.\par
\textbf{Đáp án của câu hỏi là A}
}
\end{ex}

\begin{ex} %[1D10V40]
Một sợi dây dài 65 cm. Bạn Mai cắt bớt 22 cm. Hỏi sợi dây còn lại dài bao nhiêu xăng-ti-mét?
\choice
{$42\text{ cm}$}
{$44\text{ cm}$}
{$41\text{ cm}$}
{\True $43\text{ cm}$}
\loigiai{
Chiều dài còn lại là: $65 - 22 = 43$ (cm).
\textbf{Đáp số:} $43$.\par
\textbf{Đáp án của câu hỏi là D}
}
\end{ex}

\begin{ex} %[1D10V40]
Đoạn thẳng AB dài 40 cm, đoạn thẳng CD dài 50 cm. Cả hai đoạn thẳng dài bao nhiêu xăng-ti-mét?
\choice
{$80\text{ cm}$}
{\True $90\text{ cm}$}
{$100\text{ cm}$}
{$89\text{ cm}$}
\loigiai{
Tổng độ dài hai đoạn là: $40 + 50 = 90$ (cm).
\textbf{Đáp số:} $90$.\par
\textbf{Đáp án của câu hỏi là B}
}
\end{ex}

\begin{ex} %[1D10V40]
Nếu hôm nay là Thứ Tư ngày 12 thì Thứ Tư tuần sau là ngày bao nhiêu?
\choice
{$18\text{ ngày}$}
{$20\text{ ngày}$}
{\True $19\text{ ngày}$}
{$17\text{ ngày}$}
\loigiai{
Thứ Tư tuần sau là: $12 + 7 = 19$.
\textbf{Đáp số:} $19$.\par
\textbf{Đáp án của câu hỏi là C}
}
\end{ex}

\begin{ex} %[1D10V40]
Nếu hôm nay là Thứ Bảy ngày 25 thì Thứ Bảy tuần trước là ngày bao nhiêu?
\choice
{\True $18\text{ ngày}$}
{$17\text{ ngày}$}
{$19\text{ ngày}$}
{$16\text{ ngày}$}
\loigiai{
Thứ Bảy tuần trước là: $25 - 7 = 18$.
\textbf{Đáp số:} $18$.\par
\textbf{Đáp án của câu hỏi là A}
}
\end{ex}

\textbf{Phần 2. Trắc nghiệm đúng sai}
\setcounterpageref{ex}{0}

\begin{ex} %[1D10V40]
Xét tính đúng hoặc sai của các khẳng định toán học sau đây:
\choiceTF
{Một ngày có $25$ giờ.}
{Kết quả của bài toán «Cần ít nhất bao nhiêu que tính để xếp thành 2 hình tam giác riêng biệt» là $7$.}
{\True Kết quả của bài toán «Cần ít nhất bao nhiêu que tính để xếp thành 2 hình vuông riêng biệt» là $8$.}
{\True Hình tròn có $0$ cạnh thẳng.}
\loigiai{
\begin{itemize}
    \item Ý a) Sai: Kết quả đúng là $24$, không phải $25$. Một ngày có đúng 24 giờ.
\textbf{Đáp số:} $24$.
    \item Ý b) Sai: Kết quả đúng là $6$, không phải $7$. Mỗi hình tam giác cần 3 que: $3 + 3 = 6$ (que tính).
\textbf{Đáp số:} $6$.
    \item Ý c) Đúng: Mỗi hình vuông cần 4 que: $4 + 4 = 8$ (que tính).
\textbf{Đáp số:} $8$.
    \item Ý d) Đúng: Hình tròn không có cạnh thẳng nào (có 0 cạnh thẳng).
\textbf{Đáp số:} $0$.
\end{itemize}
\begin{itemize}
    \item a) \textbf{Sai}.
    \item b) \textbf{Sai}.
    \item c) \textbf{Đúng}.
    \item d) \textbf{Đúng}.
\end{itemize}
}
\end{ex}

\begin{ex} %[1D10V40]
Xét tính đúng hoặc sai của các bài toán và phép tính sau đây:
\choiceTF
{\True Kết quả của bài toán «Một hình chữ nhật dài 25 cm, rộng 14 cm. Chiều dài hơn chiều rộng bao nhiêu xăng-ti-mét» là $11\text{ cm}$.}
{Kết quả của bài toán «Bác thợ mộc có 8 khối gỗ hình lập phương. Bác dùng 5 khối gỗ đóng giá sách. Bác còn lại bao nhiêu khối gỗ» là $2$.}
{Kết quả của bài toán «Đồng hồ chỉ 2 giờ chiều. Sau 4 giờ nữa đồng hồ chỉ mấy giờ chiều» là $7\text{ giờ}$.}
{\True Kết quả của bài toán «Cần ít nhất bao nhiêu que tính để xếp thành 2 hình vuông riêng biệt» là $8$.}
\loigiai{
\begin{itemize}
    \item Ý a) Đúng: Chiều dài hơn chiều rộng là: $25 - 14 = 11$ (cm).
\textbf{Đáp số:} $11$.
    \item Ý b) Sai: Kết quả đúng là $3$, không phải $2$. Số khối gỗ còn lại là: $8 - 5 = 3$ (khối).
\textbf{Đáp số:} $3$.
    \item Ý c) Sai: Kết quả đúng là $6$, không phải $7$. Thời điểm đó là: $2 + 4 = 6$ (giờ chiều).
\textbf{Đáp số:} $6$.
    \item Ý d) Đúng: Mỗi hình vuông cần 4 que: $4 + 4 = 8$ (que tính).
\textbf{Đáp số:} $8$.
\end{itemize}
\begin{itemize}
    \item a) \textbf{Đúng}.
    \item b) \textbf{Sai}.
    \item c) \textbf{Sai}.
    \item d) \textbf{Đúng}.
\end{itemize}
}
\end{ex}

\textbf{Phần 3. Trả lời ngắn}
\setcounterpageref{ex}{0}

\begin{ex} %[1D10V40]
Đồng hồ chỉ 11 giờ trưa. Trước đó 3 giờ là mấy giờ sáng?
\shortans{8}
\loigiai{
\textbf{Đáp số:} $8$\par
Thời điểm đó là: $11 - 3 = 8$ (giờ sáng).
\textbf{Đáp số:} $8$.
}
\end{ex}

\begin{ex} %[1D10V40]
Một tuần lễ em đi học 5 ngày. Hỏi trong 4 tuần lễ em đi học bao nhiêu ngày?
\shortans{20}
\loigiai{
\textbf{Đáp số:} $20$\par
Tổng số ngày đi học là: $5 + 5 + 5 + 5 = 20$ (ngày).
\textbf{Đáp số:} $20$.
}
\end{ex}

\begin{ex} %[1D10V40]
Một băng giấy dài 80 cm, cắt thành hai đoạn bằng nhau. Mỗi đoạn dài bao nhiêu xăng-ti-mét?
\shortans{40}
\loigiai{
\textbf{Đáp số:} $40$\par
Mỗi đoạn dài $40$ cm (vì $40 + 40 = 80$).
\textbf{Đáp số:} $40$.
}
\end{ex}

\begin{ex} %[1D10V40]
Tổng số mặt của hai khối lập phương là bao nhiêu mặt?
\shortans{12}
\loigiai{
\textbf{Đáp số:} $12$\par
Mỗi khối có 6 mặt: $6 + 6 = 12$ (mặt).
\textbf{Đáp số:} $12$.
}
\end{ex}

\begin{ex} %[1D10V40]
Tổng số đỉnh của một hình tam giác và một khối hộp chữ nhật là bao nhiêu?
\shortans{11}
\loigiai{
\textbf{Đáp số:} $11$\par
Hình tam giác có 3 đỉnh, khối hộp chữ nhật có 8 đỉnh: $3 + 8 = 11$ (đỉnh).
\textbf{Đáp số:} $11$.
}
\end{ex}

\begin{ex} %[1D10V40]
Một thước kẻ dài 20 cm, đặt que tính dài 8 cm lên thước từ vạch số 0. Đầu kia que tính chỉ vạch số mấy?
\shortans{8}
\loigiai{
\textbf{Đáp số:} $8$\par
Đầu kia chỉ đúng vạch số $8$.
\textbf{Đáp số:} $8$.
}
\end{ex}
